% Clase 20 — campo de un material magnetizado. Cada elemento dV' en r' tiene
% momento dipolar M(r') dV' y contribuye al potencial vector en r con el
% potencial de un dipolo. n es la normal saliente en la superficie S.
\begin{tikzpicture}[scale=1]

  % material
  \fill[figgray!18]
    plot[smooth cycle, tension=0.8] coordinates {(0.4,0.3) (0.9,2.3) (2.6,3.0) (4.3,2.4) (4.6,0.7) (3.2,-0.4) (1.4,-0.4)};
  \draw[figline, line width=1pt]
    plot[smooth cycle, tension=0.8] coordinates {(0.4,0.3) (0.9,2.3) (2.6,3.0) (4.3,2.4) (4.6,0.7) (3.2,-0.4) (1.4,-0.4)};
  \node[figlbl] at (1.5,0.25) {$V$};
  \node[figlbl] at (0.45,2.55) {$S$};

  % origen
  \fill[figgray] (-1.6,-1.3) circle (1.4pt);
  \node[figlblaux] at (-1.85,-1.5) {$O$};

  % elemento dV' y su momento
  \draw[figblue, line width=0.8pt, fill=figblue!20] (2.75,1.35) rectangle (3.15,1.75);
  \draw[figvecres] (2.95,1.55) -- (3.35,2.5);
  \node[figlbl, text=figred, anchor=east] at (3.1,2.4) {$\vec M\,dV'$};
  \draw[figvec] (-1.6,-1.3) -- (2.74,1.42);
  \node[figlbl] at (0.85,0.55) {$\vec r\,'$};

  % punto de observacion
  \coordinate (P) at (8.0,1.9);
  \fill[figblue] (P) circle (1.6pt);
  \draw[figvec] (-1.6,-1.3) -- (P);
  \node[figlbl] at (5.6,0.55) {$\vec r$};
  \draw[figaux, -{Latex[length=2mm,width=1.4mm]}] (3.17,1.57) -- (P);
  \node[figlbl] at (6.1,2.05) {$\vec r-\vec r\,'$};

  % normal
  \draw[figvecaux] (4.3,2.4) -- (4.95,3.0);
  \node[figlblaux] at (5.2,3.1) {$\hat n$};

\end{tikzpicture}
